% VI36-DETAILED-PROBLEM-07
\begin{problem}[VI.36.P7: Global sections of \(\mathcal O(d)\)]\label{prob:vi36-07}
Give a chartwise proof that a homogeneous polynomial \(F\in S_d\) defines a global section of \(\mathcal O(d)\), and compute the transition compatibility explicitly.
\end{problem}

\ifdefined\FullProblemDossiers
\begin{solution}
Let
\[
S=k[x_0,\ldots,x_n]
\]
and \(F\in S_d\). On the chart \(U_i\), use the local frame
\[
e_i=x_i^d
\]
of \(\mathcal O(d)\). The coefficient
\[
\frac{F}{x_i^d}
\]
has degree zero in \(S_{x_i}\), hence is a regular function on \(U_i\). Define
\[
s_i=\frac{F}{x_i^d}e_i.
\]

On \(U_i\cap U_j\),
\[
e_i
=
\left(\frac{x_i}{x_j}\right)^d e_j.
\]
Therefore
\[
s_i
=
\frac{F}{x_i^d}
\left(\frac{x_i}{x_j}\right)^d e_j
=
\frac{F}{x_j^d}e_j
=
s_j.
\]
Thus the local sections glue to a global section
\[
s_F\in\Gamma(\mathbb P^n,\mathcal O(d)).
\]
This gives the natural map
\[
S_d\longrightarrow\Gamma(\mathbb P^n,\mathcal O(d))
\]
used in the global-section theorem.
\end{solution}
\fi
